\PaperPage{45}{Appendix A.1 Verification evidence and its limits}
The formal evidence refers to a fixed source baseline, not to every
sentence in this expanded exposition. GitHub Actions run
\href{https://github.com/randogelist/Hilber-Vector-Resolvent-Theory/actions/runs/40712534224}{40712534224}
completed successfully at source commit
\begin{center}
\code{c6b95fd9e3d792270026da083bddddbc72bacb22}.
\end{center}
The preserved record reports source-identity checks, the full project
build, four exact endpoint type checks, and a strict audit of 641
selected declarations. The run's complete job log and final audit
block were independently read when that checkpoint was recorded.
The corresponding evidence files remain in the repository.

\begin{center}
\begin{tabular}{@{}ll@{}}
\toprule
Property & Recorded value\\
\midrule
Lean toolchain & \code{leanprover/lean4:v4.33.0}\\
Preserved Lean files & 118 (117 modules and the final audit)\\
Audited declarations & 641\\
Exact endpoint checks & Four, all passed\\
Active compiler target & \code{GITHUB_CI}\\
Local Lean execution for this manuscript & None\\
\bottomrule
\end{tabular}
\end{center}
The mathlib revision is pinned to
\begin{center}
\code{db584cd6d46c92f209a44c0f1c829460d327499d}.
\end{center}
The allowed axiom set is
\code{propext}, \code{Classical.choice}, and \code{Quot.sound}.
Each of the four principal endpoints has exactly this recorded set.
The audit parser rejects missing or duplicate target entries and
axioms outside the allowlist. This is evidence for the checked
declarations and their dependencies, not a replacement for reading
their types and understanding the meaning of their premises.

A separate native checkpoint records a user-reported Windows Lean
4.33.0 PASS. The original full native logs and native dependency
revision were not independently supplied. The native and GitHub
records are therefore retained separately; matching toolchain labels
alone do not establish identical dependency revisions.

The present expansion changes manuscript files only. It does not
modify the frozen Lean sources, discharge $H00$, or add a formal
implementation of the displayed complex $L^2$ operator model.
The latter is a conventional mathematical realization whose source
formalization remains a distinct extension. Historical comments in
the unchanged modules describe their preparation state before the
later PASS records; the preserved evidence identifies the subsequent
verification baseline precisely.

\PaperPage{46}{Appendix A.2 Exact endpoints and semantic assumptions}
The four exported declarations lie in the namespace
\code{RH_HilbertLimit.TwoRoutes}. With that namespace suppressed,
their checked names and types are:
\begin{quote}\small\ttfamily
nativeRH\_of\_H00\_route\_one\par
\quad : H00 $\to$ \_root\_.RiemannHypothesis\par\medskip
H00\_iff\_nativeRH\_route\_one\par
\quad : H00 $\leftrightarrow$ \_root\_.RiemannHypothesis\par\medskip
nativeRH\_of\_H00\_route\_two\par
\quad : H00 $\to$ \_root\_.RiemannHypothesis\par\medskip
H00\_iff\_nativeRH\_route\_two\par
\quad : H00 $\leftrightarrow$ \_root\_.RiemannHypothesis
\end{quote}
Here the conclusion is mathlib's literal global proposition, rather
than a newly named project-local surrogate. The premise is the
definition
\begin{center}
 \code{H00 := zeroBudgetCandidate.Admissible}.
\end{center}
Its expanded mathematical content is the zero-cancellation condition
for every actual row. Since $C_n\ge0$ unconditionally, the zero-budget
inequality $C_n\le0$ is equivalent to $C_n=0$ for every $n$.
The direct finite identities then give
\[
 H00\ \Longleftrightarrow\ (\forall n,\ C_n=0)
 \ \Longleftrightarrow\ (\forall n,\ V_n=1)
 \ \Longleftrightarrow\ (\forall N,j\le N,\ w_{N,j}\ge0).
\]
These equivalences identify the hypothesis; they do not establish it
for the actual sequence. In particular, the data pair $(0,0)$ being
finite does not make its semantic admissibility a terminating check.

The two forward constructions have a shared beginning and distinct
closures. Route~1 takes constant variation to subexponential growth
and invokes the constructed actual spectral closure. Route~2 takes
row positivity to a probability measure, obtains the actual
logarithmic-derivative germ, and invokes analytic zero exclusion.
Both reverse implications use the retained theorem
\code{H00_of_nativeRH}. Thus two equivalence names do not assert two
independent proofs of every supporting lemma.

For interpretation of other audited declarations, their hypotheses
remain decisive. A theorem about a supplied
\code{BKQRContractiveOrbit} proves consequences of that structure;
it does not construct the structure for the actual Widder family.
Similarly, a theorem whose premise is
\code{ActualApproximationConvergence} does not remove that premise.
The audit confirms proof terms for the stated propositions, while
the manuscript's claim correspondence tracks which objects and
assumptions those propositions describe.

\PaperPage{47}{Appendix B.1 What a finite positive prefix can establish}
Finite rows provide necessary tests for an infinite moment sequence.
They do not, for an unrestricted sequence, force positivity of all
later rows. The following elementary example makes the quantifiers
explicit.

\begin{proposition}\label{app:prefix}
For every integer $K\ge0$, there is a real normalized sequence whose
Bernstein rows of degrees $0,\ldots,K$ are nonnegative but whose row
of degree $K+1$ has a negative entry.
\end{proposition}
\begin{proof}
Set
\[
 u_0=1,\qquad u_1=\cdots=u_K=0,\qquad u_{K+1}=-1,
\]
and, for definiteness, set all later terms to zero. When $K=0$,
the middle condition is empty. A row of degree $N$ only uses the
terms $u_0,\ldots,u_N$. Hence for $N\le K$ its weights agree with
those of the point mass $\delta_0$:
\[
 w_{N,0}=1,\qquad w_{N,j}=0\quad(1\le j\le N).
\]
Every such row is nonnegative and has variation one. In the next
row, the last weight is
\[
 w_{K+1,K+1}=u_{K+1}=-1<0,
\]
which proves the assertion.
\end{proof}
The example also has $C_n=0$ for $n<K$, because those transitions
join two rows of variation one. The next variation is strictly
greater than one: row mass is still one but at least one weight is
negative. The identity $C_K=(K+1)(V_{K+1}-V_K)$ therefore detects a
strictly positive cancellation at the first unchecked transition.

This obstruction is about arbitrary extensions of finite scalar data.
It does not prove that no finite certificate can ever establish a
universal property of a fixed arithmetic sequence. A successful
certificate could include additional structure, an induction
invariant, or a theorem controlling every remaining row. What the
example excludes is the inference from a positive prefix alone,
without such a propagation argument.

For the actual $\xi$ sequence, a negative exact weight would refute
$H00$, whereas any finite list of nonnegative weights is consistent
with both the desired universal assertion and failure farther out.
Numerical approximations introduce a further distinction: a displayed
small number is not an exact sign proof unless its error is controlled.
The present endpoints start from semantic $H00$ and do not supply a
finite algorithm deciding it. This is why the finite-checker program
in Part~5 remains an additional theorem target.

\PaperPage{48}{Appendix B.2 Quantitative approximation of the resolvent}
The moment construction also gives useful quantitative estimates once
the representing measure is available. Define
$\sigma_N=\sum_{j=0}^N w_{N,j}\delta_{j/N}$ for $N\ge1$.
Using the moment formula in each Bernstein weight gives
\[
 \int f\,d\sigma_N
 =\int_0^1 B_Nf(q)\,d\sigma(q),\qquad
 B_Nf(q)=\sum_{j=0}^N\binom Nj q^j(1-q)^{N-j}f(j/N).
\]
For fixed $q$, this is the expectation of $f(X/N)$ with
$X$ binomial of parameters $N,q$. If $f$ is complex-valued and
$L$-Lipschitz on $[0,1]$, Cauchy--Schwarz and the binomial variance give
\[
 |B_Nf(q)-f(q)|
 \le L\,\mathbb E|X/N-q|
 \le L\sqrt{\frac{q(1-q)}N}\le\frac{L}{2\sqrt N}.
\]
Integrating proves the same upper bound for
$|\int f\,d\sigma_N-\int f\,d\sigma|$. For continuous $f$,
uniform approximation by Lipschitz functions yields weak convergence
$\sigma_N\Rightarrow\sigma$. This supplies an ordinary limiting
interpretation of the measure after its construction.

\begin{proposition}\label{app:resolvent-rate}
Let $K$ be a nonempty compact subset of
$\mathbb C\setminus(-\infty,-1/\kappa]$. Put
\[
 d=\min_{z\in K,\,q\in[0,1]}|1+\kappa qz|>0,
 \qquad M=\max_{z\in K}|z|.
\]
For $S_{\kappa,N}(z)=\kappa\int(1+\kappa qz)^{-1}\,d\sigma_N(q)$,
\[
 \sup_{z\in K}|S_{\kappa,N}(z)-S_\kappa(z)|
 \le\frac{\kappa^2M}{2d^2\sqrt N}.
\]
\end{proposition}
\begin{proof}
For $f_z(q)=\kappa/(1+\kappa qz)$, the derivative with respect to
the real variable $q$ has modulus at most $\kappa^2M/d^2$.
The fundamental theorem of calculus makes this a uniform Lipschitz
bound on the interval. Apply the preceding estimate for each $z$.
\end{proof}
Thus positive finite rows approximate the scalar resolvent uniformly
on compact subsets of its common domain. The constants deteriorate
near the slit because denominators can become small. This theorem
does not identify the distinct weighted Widder coefficients: it
concerns the measure resolvent built from the already specified
moments. The estimate is a paper-level consequence of that measure
representation, not an additional claim about the recorded Lean audit.

\PaperPage{49}{References}
\begingroup
\renewcommand{\section}[2]{}%
\begin{thebibliography}{9}
\setlength{\itemsep}{9pt}
\bibitem{conway}
J.~B. Conway, \emph{A Course in Functional Analysis}, second edition,
Graduate Texts in Mathematics~96, Springer-Verlag, New York, 1990.
\href{https://doi.org/10.1007/978-1-4757-4383-8}{doi:10.1007/978-1-4757-4383-8}.

\bibitem{hausdorff}
F.~Hausdorff, Summationsmethoden und Momentfolgen.~I,
\emph{Mathematische Zeitschrift} \textbf{9} (1921), 74--109.
\href{https://doi.org/10.1007/BF01378337}{doi:10.1007/BF01378337}.

\bibitem{lean4}
L.~de Moura and S.~Ullrich, The Lean 4 theorem prover and programming
language, in \emph{Automated Deduction---CADE 28}, Lecture Notes in
Computer Science~12699, Springer, 2021, pp.~625--635.
\href{https://doi.org/10.1007/978-3-030-79876-5_37}{doi:10.1007/978-3-030-79876-5\_37}.

\bibitem{mathlib}
The mathlib Community, The Lean mathematical library,
in \emph{Proceedings of the 9th ACM SIGPLAN International Conference
on Certified Programs and Proofs (CPP 2020)}, ACM, 2020, pp.~367--381.
\href{https://doi.org/10.1145/3372885.3373824}{doi:10.1145/3372885.3373824}.

\bibitem{rudin}
W.~Rudin, \emph{Real and Complex Analysis}, third edition,
McGraw-Hill, New York, 1987. ISBN~978-0-07-054234-1.

\bibitem{titchmarsh}
E.~C. Titchmarsh, \emph{The Theory of the Riemann Zeta-function},
second edition, revised by D.~R. Heath-Brown,
Clarendon Press, Oxford, 1986. ISBN~0-19-853369-1.

\bibitem{source}
\emph{Hilber-Vector-Resolvent-Theory}, project source repository,
frozen construction package
\code{RH_H00_TWO_ROUTES_LEAN_v0_2_3_construction}.
Audited source commit \code{c6b95fd9e3d792270026da083bddddbc72bacb22}.
\href{https://github.com/randogelist/Hilber-Vector-Resolvent-Theory}{Repository and manuscript sources};
\href{https://github.com/randogelist/Hilber-Vector-Resolvent-Theory/actions/runs/40712534224}{GitHub CI run 40712534224}.
Verification records dated 28 October 2026.
\end{thebibliography}
\endgroup
\medskip
The classical references provide background for the Hilbert-space,
moment, and analytic constructions. The source reference identifies
the exact project statements and their proof dependencies. The Lean
and mathlib publications describe the proof infrastructure; the
specific dependency versions used by this project are the pins in
Appendix~A. No new bibliography entry is intended to replace a
source-level audit of the actual conditional endpoints.

\PaperPage{50}{Appendix C. Source concordance and reproducibility}
The table locates the main mathematical stages. Paths in the middle
column are relative to \code{RH_HilbertLimit/TwoRoutes/}, except where
an explicit parent path is shown. The source baseline is the commit
recorded in Appendix~A.

\begin{center}\small
\begin{tabular}{@{}p{0.24\linewidth}p{0.45\linewidth}p{0.23\linewidth}@{}}
\toprule
Manuscript topic & Source module or record & Role\\
\midrule
Hilbert defects & \code{Derivations/BKQR/}\newline\code{DefectGramSquare.lean} & Abstract real Gram interface\\[5pt]
$H00$ and row positivity & \code{RouteOne.lean};\newline\code{ActualPositivePrefix.lean} & Actual conditional identities\\[5pt]
Monomial limit & \code{BernsteinMomentLimit.lean} & General real sequences\\[5pt]
Positive measure & \code{HausdorffMeasure.lean} & Constructed from positivity\\[5pt]
Kernel holomorphy & \code{ResolventAnalytic.lean} & Local denominator control\\[5pt]
Actual germ & \code{MeasureResolvent.lean} & Exact derivative identification\\[5pt]
Continuation & \code{ActualResolventContinuation.lean} & Analytic defect argument\\[5pt]
Pole closure & \code{ActualHerglotzClosure.lean} & Sign implies actual RH\\[5pt]
Weighted frontier & \code{RouteTwoContractive.lean};\newline\code{ActualLimitBridge.lean} & Bounds and explicit open bridge\\[5pt]
Endpoint types & \code{ci/EndpointTypes.lean} & Four exact type checks\\[5pt]
Audit evidence & \code{evidence/}\newline\code{GITHUB_CI_AXIOM_AUDIT.json} & 641 declaration results\\
\bottomrule
\end{tabular}
\end{center}

The spectral route also uses the retained
\code{RH_HilbertLimit/} modules. These include
\code{AnalyticClosure.lean} and the actual calibration and extraction
infrastructure. The finer
introduction source map under \code{Paper-Candidates/} records the
difference between the minimal real Gram interface and the concrete
complex Hilbert-space realization used in this exposition.

To rebuild this manuscript, run \code{python build.py} in its source
directory with Python, \code{pdflatex}, and Poppler's \code{pdfinfo}
available. The builder invokes only LaTeX, checks the expected
subsection-to-page map, rejects overflow and unresolved references,
and verifies the fifty-page output. Lean compilation remains assigned
to the existing GitHub workflow. Rebuilding the PDF is not a new
formal-verification checkpoint.

The next mathematical extensions have separate acceptance criteria:
construct the concrete operator model in Lean; prove the exact Widder
identification if the longer route is pursued; and establish a
universal admissibility argument before asserting an accepted
certificate or unconditional RH. Each would change the theorem
coverage and should carry its own source and verification record.
